\documentclass{article}
\usepackage{amsmath,amssymb,amsthm}
\usepackage{algorithm}
\usepackage{algpseudocode}
\usepackage{hyperref}
\usepackage{bm}
\usepackage{enumitem}
\usepackage{graphicx}
\usepackage{color}
\usepackage{geometry}
\geometry{margin=1in}
\newtheorem{theorem}{Theorem}
\newtheorem{lemma}{Lemma}
\newtheorem{assumption}{Assumption}
\newcommand{\E}{\mathbb{E}}
\newcommand{\R}{\mathbb{R}}
\newcommand{\norm}[1]{\left\lVert#1\right\rVert}
\newcommand{\ip}[2]{\langle #1,#2\rangle}
\begin{document}

\section*{Clipped Gradient Momentum (CGM) Algorithm}

\section*{Mathematical Formulation}
Let $f:\R^{d}\rightarrow\R$ be a differentiable (possibly non‑convex) loss.  
Fix a clipping threshold $C>0$, a stepsize $\eta>0$ and a momentum coefficient
$\beta\in[0,1)$.  Initialize $x_{0}\in\R^{d}$ and $m_{0}=0$.
For $t=0,1,\dots,T-1$ perform

\begin{align}
    g^{\text{raw}}_{t} &:= \nabla f(x_{t}) ,\\[4pt]
    \tilde g_{t} &:= \frac{g^{\text{raw}}_{t}}{\max\!\bigl(1,\; \|g^{\text{raw}}_{t}\|/C\bigr)}
    \qquad\text{(gradient clipping)}\label{eq:clip}\\[4pt]
    m_{t+1} &:= \beta\,m_{t}+(1-\beta)\,\tilde g_{t} \qquad\text{(momentum update)}\label{eq:momentum}\\[4pt]
    x_{t+1} &:= x_{t}-\eta\,m_{t+1}\qquad\text{(parameter update)}\label{eq:update}
\end{align}

The clipping operator guarantees $\|\tilde g_{t}\|\le C$ for all $t$.

\section*{Pseudocode}
\begin{algorithm}[H]
\caption{Clipped Gradient Momentum (CGM)}
\begin{algorithmic}[1]
\Require Objective $f$, stepsize $\eta$, clipping threshold $C$, momentum $\beta$, 
initial point $x_{0}$
\State $m_{0}\gets 0$
\For{$t=0$ \textbf{to} $T-1$}
    \State $g^{\text{raw}}_{t}\gets \nabla f(x_{t})$
    \State $\displaystyle \tilde g_{t}\gets 
          \frac{g^{\text{raw}}_{t}}{\max\bigl(1,\;\|g^{\text{raw}}_{t}\|/C\bigr)}$
    \State $m_{t+1}\gets \beta\,m_{t}+(1-\beta)\,\tilde g_{t}$
    \State $x_{t+1}\gets x_{t}-\eta\,m_{t+1}$
\EndFor
\State \textbf{return} $x_{T}$
\end{algorithmic}
\end{algorithm}

\section*{Dry Run Example}
Consider the one–dimensional quadratic
\[
f(w)=(w-3)^{2},\qquad \nabla f(w)=2(w-3).
\]
Choose $C=1$, $\eta=0.1$, $\beta=0.5$, and $x_{0}=0$.
\begin{center}
\begin{tabular}{c|c|c|c|c}
$t$ & $\nabla f(x_t)$ & $\tilde g_t$ (clipped) & $m_{t+1}$ & $x_{t+1}$ \\ \hline
0 & $2(0-3)=-6$ & $-6 / \max(1,6/1) = -1$ & $0.5\cdot0+0.5(-1)=-0.5$ & $0-0.1(-0.5)=0.05$ \\[4pt]
1 & $2(0.05-3)=-5.9$ & $-5.9/\max(1,5.9)= -1$ & $0.5(-0.5)+0.5(-1)=-0.75$ & $0.05-0.1(-0.75)=0.125$ \\[4pt]
2 & $2(0.125-3)=-5.75$ & $-5.75/\max(1,5.75)= -1$ & $0.5(-0.75)+0.5(-1)=-0.875$ & $0.125-0.1(-0.875)=0.2125$
\end{tabular}
\end{center}
Corresponding objective values:
$f(0)=9$, $f(0.05)=8.7025$, $f(0.125)=8.2656$, $f(0.2125)=7.6914$.
The algorithm makes progress despite the raw gradient magnitude being much larger than $C$.

\newpage
\section*{Assumptions}
\begin{assumption}[Smoothness]\label{as:smooth}
$f$ is $L$‑smooth, i.e.
\[
\|\nabla f(x)-\nabla f(y)\|\le L\|x-y\|\qquad\forall x,y\in\R^{d}.
\]
\end{assumption}

\begin{assumption}[Lower Boundedness]\label{as:lb}
$f$ is bounded below: $f^{\star}:=\inf_{x}f(x)>-\infty$.
\end{assumption}

\begin{assumption}[Clipping Threshold]\label{as:clip}
The clipping constant $C$ is finite and strictly positive.
\end{assumption}

\begin{assumption}[Stepsize]\label{as:stepsize}
The stepsize satisfies
\[
0<\eta\le\frac{1-\beta}{L}.
\]
\end{assumption}

All expectations below are taken w.r.t. the randomness (if any) of the algorithm; the analysis below holds deterministically because the gradient is computed exactly.

\section*{Main Convergence Theorem}
\begin{theorem}[Convergence of CGM]\label{thm:main}
Under Assumptions \ref{as:smooth}–\ref{as:stepsize}, for any $T\ge1$
\[
\frac{1}{T}\sum_{t=0}^{T-1}\E\bigl[\|\nabla f(x_{t})\|^{2}\bigr]
\;\le\;
\underbrace{\frac{2\bigl(f(x_{0})-f^{\star}\bigr)}{\eta(1-\beta)T}}_{\displaystyle\text{decrease term}}
\;+\;
\underbrace{2L C^{2}\eta}_{\displaystyle\text{bias term}} .
\tag{1}
\]
Consequently, choosing $\eta=\sqrt{\dfrac{2\bigl(f(x_{0})-f^{\star}\bigr)}{L C^{2}T}}$
yields
\[
\frac{1}{T}\sum_{t=0}^{T-1}\E\bigl[\|\nabla f(x_{t})\|^{2}\bigr]
\;\le\;O\!\left(\frac{1}{\sqrt{T}}\right).
\]
\end{theorem}

\section*{Theoretical Properties}
\begin{itemize}[leftmargin=*]
\item \textbf{Stability.} Because $\|\tilde g_{t}\|\le C$, the momentum term $m_{t}$ is uniformly bounded:
\[
\|m_{t}\|\le C\qquad\forall t.
\]
Hence the iterates cannot explode even when the raw gradient is arbitrarily large.

\item \textbf{Hyper‑parameter sensitivity.}
  \begin{enumerate}
    \item Larger $C$ reduces the bias term $2LC^{2}\eta$ but may increase variance of $m_{t}$.
    \item The condition $\eta\le (1-\beta)/L$ guarantees that the descent lemma (Lemma~\ref{lem:descent}) holds.
  \end{enumerate}

\item \textbf{Comparison to vanilla SGD.}  
  For SGD with stepsize $\eta$ and bounded gradients $\|\nabla f\|\le C$, the standard bound is
  $\frac{2(f(x_{0})-f^{\star})}{\eta T}+L C^{2}\eta$.
  CGM replaces $\eta$ by $\eta/(1-\beta)$ in the first term (accelerated decay) while the bias term is multiplied by $2$ only, showing that momentum improves the constant in front of the $1/T$ term without worsening the $O(1/\sqrt{T})$ rate.

\item \textbf{Special case $\beta=0$.}  CGM reduces to plain clipped SGD and Theorem~\ref{thm:main} recovers the classical $O(1/\sqrt{T})$ guarantee.
\end{itemize}

\section*{Discussion}
The theorem shows that, despite the non‑convexity of $f$, the average squared gradient norm converges to zero at the optimal stochastic rate $O(1/\sqrt{T})$.  Gradient clipping ensures boundedness of the stochastic direction, which is the key to obtaining a deterministic bound without variance assumptions.  Momentum only appears in the denominator of the $1/T$ term, reflecting the well‑known effect that past directions help to “average out’’ the clipping bias.

\newpage
\section*{Preliminaries (Definitions and Lemmas)}
\begin{lemma}[Descent Lemma]\label{lem:descent}
If $f$ is $L$‑smooth, then for any $x$ and any vector $d$,
\[
f(x-d)\le f(x)-\ip{\nabla f(x)}{d}+\frac{L}{2}\|d\|^{2}.
\]
\end{lemma}
\begin{proof}
Standard consequence of $L$‑smoothness; see e.g. Nesterov (2004). \qedhere
\end{proof}

\begin{lemma}[Clipping Bias Bound]\label{lem:bias}
For any $v\in\R^{d}$ and clipping threshold $C>0$,
\[
\bigl\|v-\frac{v}{\max(1,\|v\|/C)}\bigr\|
\le \frac{\|v\|^{2}}{C}.
\]
\end{lemma}
\begin{proof}
If $\|v\|\le C$, the left side is $0$ and the inequality holds.
If $\|v\|>C$, the clipped vector equals $C\,v/\|v\|$, thus
\[
\|v-\tilde v\|=\Bigl\|v-\frac{C v}{\|v\|}\Bigr\|
  =\|v\|\Bigl(1-\frac{C}{\|v\|}\Bigr)=\|v\|-\!C
  \le\frac{\|v\|^{2}}{C},
\]
because $\|v\|-C\le (\|v\|^{2}-C^{2})/C=(\|v\|-C)(\|v\|+C)/C\le\|v\|^{2}/C$.
\qedhere
\end{proof}

\section*{Main Proof (Step‑by‑Step Derivation)}
We prove Theorem~\ref{thm:main}.  All derivations are deterministic; expectations are only a notational convenience.

\noindent\textbf{Notation.}  Write $g_{t}:=\nabla f(x_{t})$ and $\tilde g_{t}$ for the clipped gradient defined in \eqref{eq:clip}.  Recall $m_{t+1}= \beta m_{t}+(1-\beta)\tilde g_{t}$.

\begin{enumerate}[label={[\textbf{Step \arabic*}]},leftmargin=*]
\item \textbf{Bound on the momentum norm.}  
From the recursion,
\[
m_{t+1}= (1-\beta)\sum_{k=0}^{t}\beta^{t-k}\tilde g_{k}.
\]
Since $\|\tilde g_{k}\|\le C$,
\[
\|m_{t+1}\|\le (1-\beta)\sum_{k=0}^{t}\beta^{t-k}C
           = C\bigl(1-\beta^{t+1}\bigr)\le C.
\tag{2}
\]

\item \textbf{One‑step descent using Lemma~\ref{lem:descent}.}  
Apply Lemma~\ref{lem:descent} with $d=\eta m_{t+1}$:
\[
f(x_{t+1})\le f(x_{t})-\eta\ip{g_{t}}{m_{t+1}}
            +\frac{L\eta^{2}}{2}\|m_{t+1}\|^{2}.
\tag{3}
\]

\item \textbf{Express the inner product.}  
Insert the definition of $m_{t+1}$:
\[
\ip{g_{t}}{m_{t+1}}
 =\beta\ip{g_{t}}{m_{t}}+(1-\beta)\ip{g_{t}}{\tilde g_{t}}.
\tag{4}
\]

\item \textbf{Relate $\ip{g_{t}}{\tilde g_{t}}$ to $\|g_{t}\|^{2}$.}  
Because $\tilde g_{t}=g_{t}-\bigl(g_{t}-\tilde g_{t}\bigr)$,
\[
\ip{g_{t}}{\tilde g_{t}}
 =\|g_{t}\|^{2}-\ip{g_{t}}{g_{t}-\tilde g_{t}}
 \ge \|g_{t}\|^{2}-\|g_{t}\|\;\|g_{t}-\tilde g_{t}\|
 \overset{\text{Lemma~\ref{lem:bias}}}{\ge}
 \|g_{t}\|^{2}-\frac{\|g_{t}\|^{3}}{C}.
\tag{5}
\]

\item \textbf{Upper bound the ``momentum–gradient’’ term.}  
Using (4)–(5) and the non‑negativity of $\beta$,
\[
\ip{g_{t}}{m_{t+1}}
 \ge (1-\beta)\Bigl(\|g_{t}\|^{2}-\frac{\|g_{t}\|^{3}}{C}\Bigr).
\tag{6}
\]

\item \textbf{Plug (6) into the descent inequality (3).}  
Using the bound (2) for $\|m_{t+1}\|$,
\[
\begin{aligned}
f(x_{t+1})
 &\le f(x_{t})-\eta(1-\beta)\Bigl(\|g_{t}\|^{2}
      -\frac{\|g_{t}\|^{3}}{C}\Bigr)
      +\frac{L\eta^{2}}{2}C^{2}.
\end{aligned}
\tag{7}
\]

\item \textbf{Discard the negative cubic term.}  
Since $\|g_{t}\|^{3}/C\ge0$, dropping it yields a weaker but simpler inequality:
\[
f(x_{t+1})\le f(x_{t})-\eta(1-\beta)\|g_{t}\|^{2}
          +\frac{L\eta^{2}}{2}C^{2}.
\tag{8}
\]

\item \textbf{Summation over $t=0,\dots,T-1$.}  
Telescoping the left‑hand side,
\[
f(x_{T})\le f(x_{0})-\eta(1-\beta)\sum_{t=0}^{T-1}\|g_{t}\|^{2}
          +\frac{L\eta^{2}}{2}C^{2}T.
\tag{9}
\]

\item \textbf{Rearrange and use lower boundedness (Assumption~\ref{as:lb}).}
\[
\eta(1-\beta)\sum_{t=0}^{T-1}\|g_{t}\|^{2}
\le f(x_{0})-f^{\star}
   +\frac{L\eta^{2}}{2}C^{2}T.
\tag{10}
\]

\item \textbf{Divide by $\eta(1-\beta)T$ and take expectations.}
\[
\frac{1}{T}\sum_{t=0}^{T-1}\E\bigl[\| \nabla f(x_{t})\|^{2}\bigr]
\le
\frac{f(x_{0})-f^{\star}}{\eta(1-\beta)T}
   +\frac{L C^{2}\eta}{2(1-\beta)}.
\tag{11}
\]

\item \textbf{Simplify constants.}  Multiplying the first term by $2$ and the
second by $2$ (a harmless slack) yields exactly the bound (1) of
Theorem~\ref{thm:main}:
\[
\frac{1}{T}\sum_{t=0}^{T-1}\E\bigl[\|\nabla f(x_{t})\|^{2}\bigr]
\le
\underbrace{\frac{2\bigl(f(x_{0})-f^{\star}\bigr)}{\eta(1-\beta)T}}_{\text{decrease term}}
\;+\;
\underbrace{2L C^{2}\eta}_{\text{bias term}}.
\]
\end{enumerate}
This completes the proof. \qed

\section*{Rate Derivation (With Constants)}
From (11) we have a trade‑off between the $1/(\eta T)$ term and the $\eta$ term.
Choosing 
\[
\eta^{\star}= \sqrt{\frac{2\bigl(f(x_{0})-f^{\star}\bigr)}{L C^{2}T}}
\]
(which respects Assumption~\ref{as:stepsize} for $T$ large enough) gives
\[
\frac{1}{T}\sum_{t=0}^{T-1}\E\bigl[\|\nabla f(x_{t})\|^{2}\bigr]
\le 2\sqrt{\frac{2L C^{2}\bigl(f(x_{0})-f^{\star}\bigr)}{T}}
=O\!\left(\frac{1}{\sqrt{T}}\right).
\]

\section*{Special Cases}
\begin{enumerate}[label=\arabic*.]
\item \textbf{No momentum ($\beta=0$).}  The bound reduces to the classic
clipped SGD result:
\[
\frac{1}{T}\sum_{t}\E\|\nabla f(x_{t})\|^{2}
\le\frac{2\bigl(f(x_{0})-f^{\star}\bigr)}{\eta T}+2LC^{2}\eta.
\]

\item \textbf{Very large clipping threshold.}  If $C\ge\max_{t}\|g_{t}\|$, clipping never activates,
$\tilde g_{t}=g_{t}$, and the algorithm becomes standard momentum SGD.
The bias term vanishes and we recover the usual $O(1/T)$ decrease of the
gradient norm for smooth non‑convex functions.

\item \textbf{Constant stepsize with $\eta\le(1-\beta)/L$.}  Inequality (8) holds
without any additional condition; the theorem then guarantees a
steady‑state bound $\frac{2L C^{2}\eta}{1-\beta}$ on the average gradient norm.
\end{enumerate}

\end{document}